\chapter{File Organization and Disk-Block Allocation}\label{ch:filesystems}
\chapterguide{Chapter 2: file I/O and Chapter 9: names, inodes, and file-system objects.}{separate application record formats from file-system block placement, compare text and binary records, and trace indexed allocation in Lab 10.}

\section{A file is a byte sequence; applications add meaning}
To the operating system, a regular file can be treated as a sequence of bytes plus metadata. Applications decide how those bytes represent records such as students, employees, images, or database pages.

This gives us two layers to keep separate:
\begin{itemize}
  \item \textbf{application organization:} how fields and records are encoded inside the file;
  \item \textbf{file-system allocation:} how logical file blocks are mapped onto storage.
\end{itemize}

\section{Text and binary records}
Text formats convert values into readable characters. Binary formats can store compact machine-oriented representations.

\begin{utptable}
\renewcommand{\arraystretch}{1.18}
\begin{tabularx}{\linewidth}{@{}>{\bfseries\leavevmode\color{UTPNavy}}m{0.19\linewidth}
  >{\raggedright\arraybackslash}X
  >{\raggedright\arraybackslash}X@{}}
\toprule
\rowcolor{UTPPaleBlue!92}
\utpth{Format} & \utpth{Strength} & \utpth{Caution} \\
\midrule
Text & readable with ordinary tools, easy to exchange & parsing and formatting add overhead \\
Raw C struct & compact and easy for a controlled exercise & padding, endianness, and type sizes can reduce portability \\
\bottomrule
\end{tabularx}
\end{utptable}

\begin{pitfall}
``Binary'' does not automatically mean ``better.'' A durable binary format normally defines its byte layout explicitly instead of simply dumping a compiler-dependent in-memory structure.
\end{pitfall}

\section{Sequential access and sequential search}
A small record file can be searched by reading one record after another until the key is found. This is easy to implement and understand, but the worst case may inspect every record.

Larger systems introduce indexes so a key can point more directly to the needed location. The idea is similar to a book index: extra structure is maintained so lookup does not require scanning every page.

\section{Logical blocks and physical placement}
A file system commonly views a file as logical blocks numbered 0, 1, 2, and so on. The allocation method determines where those logical blocks are stored physically.

\section{Contiguous allocation}
All file blocks occupy consecutive storage blocks. Sequential and random access are simple, but file growth can be difficult when the following physical blocks are already occupied.

\section{Linked allocation}
Each allocated block identifies the next block in the file. Files can grow without requiring adjacent space, but following a long chain is inefficient for direct/random access and damaged pointers are a concern.

\section{Indexed allocation}
A separate index block stores references to the file's data blocks. Logical block 0 can be found through the first index entry, logical block 1 through the second, and so on. The data blocks themselves do not need to be adjacent.

\begin{utpfig}[Indexed allocation]{Indexed allocation. Index block 7 lists the file's data blocks in logical order, so the blocks themselves need not be adjacent.}{fig:indexed-allocation}
\begin{tikzpicture}[
  b/.style={draw=UTPBlue,fill=UTPPaleBlue,rounded corners=1mm,minimum width=2.2cm,minimum height=0.7cm,font=\small\bfseries\color{UTPNavy}},
  d/.style={draw=UTPTeal,fill=UTPPaleTeal,rounded corners=1mm,minimum width=1.7cm,minimum height=0.65cm,font=\scriptsize\color{UTPNavy}},
  arr/.style={-{Stealth[length=2mm]},draw=UTPMuted,thick}
]
\node[b] (idx) {Index block 7};
\node[d,right=2.2cm of idx,yshift=1.0cm] (d1) {block 12};
\node[d,right=2.2cm of idx] (d2) {block 19};
\node[d,right=2.2cm of idx,yshift=-1.0cm] (d3) {block 31};
\draw[arr] (idx)--(d1); \draw[arr] (idx)--(d2); \draw[arr] (idx)--(d3);
\end{tikzpicture}
\end{utpfig}

\begin{utptable}
\renewcommand{\arraystretch}{1.18}
\begin{tabularx}{\linewidth}{@{}>{\bfseries\leavevmode\color{UTPNavy}}m{0.18\linewidth}
  >{\raggedright\arraybackslash}m{0.40\linewidth}
  >{\raggedright\arraybackslash}X@{}}
\toprule
\rowcolor{UTPPaleBlue!92}
\utpth{Method} & \utpth{Main advantage} & \utpth{Main limitation} \\
\midrule
Contiguous & simple, efficient direct access & growth and external fragmentation \\
Linked & easy growth without contiguous space & poor direct access \\
Indexed & direct access without contiguous data blocks & index-space overhead \\
\bottomrule
\end{tabularx}
\end{utptable}

\section{Free-space management}
A file system also needs to know which storage blocks are free. A bitmap is one common idea: each bit records whether one block is free or allocated. Other structures can use lists or grouped free blocks.

The Lab 10 simulator uses a simple integer array as an easy-to-see allocation map. It is a model, not a complete disk file system.

\section{Validate first, commit second}
Suppose a request needs one index block and three data blocks. If the program marks the first two blocks as used and only later discovers that the third is unavailable, it has partially changed state.

A safer introductory pattern is:
\begin{enumerate}
  \item validate the complete request without modifying shared state;
  \item if every block is valid and free, commit all allocation changes;
  \item otherwise, leave the map unchanged.
\end{enumerate}
This transaction-like reasoning appears throughout systems software.

\section{Durability and crash consistency}
Real file systems must survive interruptions during metadata updates. Techniques such as journaling, copy-on-write, checksums, and careful write ordering help prevent inconsistent states after a crash. These mechanisms are beyond the Lab 10 implementation, but the validate-then-commit habit is an accessible first step toward the same design concern.

\section{Keep the layers separate}
When an application stores a student record, it decides which bytes represent the ID, name, and GPA. When the file system stores the file, it decides how file offsets map to blocks. One file can change its physical allocation without changing its application record format.

\begin{quickcheck}
\begin{enumerate}
  \item Which layer decides how an employee record is encoded?
  \item Which allocation method uses a dedicated block of pointers?
  \item Why should a multi-block allocation request be fully checked before state is modified?
\end{enumerate}
\end{quickcheck}

\labsection{10}{File records and indexed block allocation}{sec:lab10}
\begin{labproject}{Separate application record format from file-system block placement}

\textbf{Coverage.} Chapter~\ref{ch:filesystems}: record encoding vs. file-system allocation.\\
\textbf{Goal.} Create text and binary record files, then simulate an index block that points to data blocks.\\
\textbf{Explain.} Explain which decisions belong to the application and which belong to the file system.

\section*{Source code}
\labsourcecard{\texttt{text\_records.c}}{Write readable employee records, reopen the file, and search by ID.}{\texttt{fprintf}/\texttt{fscanf} define an application-level text format.}
\lstinputlisting[style=utpcnumbered]{all_experiments/Lab10/text_records.c}

\labsourcecard{\texttt{binary\_records.c}}{Write fixed C structs as binary data and search them after reopening the file.}{The exercise is compact but the raw struct layout is compiler/platform dependent.}
\lstinputlisting[style=utpcnumbered]{all_experiments/Lab10/binary_records.c}

\labsourcecard{\texttt{indexed\_allocation.c}}{Represent a small disk-block map and validate an index block plus data blocks before committing allocation.}{The program checks the whole request first, so a rejected file does not partially modify the map.}
\lstinputlisting[style=utpcnumbered]{all_experiments/Lab10/indexed_allocation.c}

\section*{Task 1 --- Create text and binary record files}
\begin{lstlisting}[style=utpbash]
cd all_experiments/Lab10
gcc -Wall -Wextra -std=c11 text_records.c -o text_records
gcc -Wall -Wextra -std=c11 binary_records.c -o binary_records
./text_records
./binary_records
ls -l employees.txt students.bin
cat employees.txt
od -Ax -tx1z students.bin | head
\end{lstlisting}
Compare readability, file size, and portability in a short table.

\section*{Task 2 --- Identify the layer responsible}
For each statement, label it \textbf{application} or \textbf{file system}: ``ID is stored before name'', ``file offset maps to a disk block'', ``salary is printed with two decimals'', ``an index block points to data blocks''.

\section*{Task 3 --- Indexed allocation}
\begin{lstlisting}[style=utpbash]
gcc -Wall -Wextra -std=c11 indexed_allocation.c -o indexed_allocation
./indexed_allocation
\end{lstlisting}
Explain why File B is rejected and why the block map should remain unchanged after that rejection.

\section*{Task 4 --- Controlled edit}
Open \texttt{indexed\_allocation.c} in Vim. Change File B so its data-block list no longer overlaps block 8. Predict whether it will now allocate and whether that affects File C. Recompile and verify.

\begin{utptable}
\renewcommand{\arraystretch}{1.18}
\begin{tabularx}{\linewidth}{@{}>{\bfseries\leavevmode\color{UTPNavy}}m{0.42\linewidth}
  >{\raggedright\arraybackslash}X
  >{\raggedright\arraybackslash}X@{}}
\toprule
\rowcolor{UTPPaleBlue!92}
\utpth{Question} & \utpth{Text records} & \utpth{Binary records} \\
\midrule
Human-readable directly? & & \\
Easy to inspect with normal text tools? & & \\
Raw representation portable across systems? & & \\
\bottomrule
\end{tabularx}
\end{utptable}

\begin{quickcheck}
\begin{enumerate}
  \item Which decisions belong to the application record format rather than the file system?
  \item Why is writing a raw C structure not a portable long-term file format?
  \item What role does an index block play in indexed allocation?
\end{enumerate}
\end{quickcheck}

\begin{verification}
\begin{itemize}
  \item The text and binary record files are compared on readability, size and portability.
  \item Each of the four items is assigned to the layer responsible for it.
  \item File B is rejected and the block map is unchanged after the rejection.
  \item The before/after block map from your File B edit is explained.
\end{itemize}
\end{verification}

\begin{labdeliverable}
Submit the text/binary comparison, the four layer labels, the original indexed-allocation output, and the before/after block-map result from your File B edit.
\end{labdeliverable}
\end{labproject}

